\documentclass[10pt]{article}
\usepackage{fontspec}
\setmainfont{lmroman10-regular.otf}[BoldFont=lmroman10-bold.otf,ItalicFont=lmroman10-italic.otf,BoldItalicFont=lmroman10-bolditalic.otf]
\setmonofont{lmmono10-regular.otf}
\usepackage[a4paper,margin=16mm,top=14mm,bottom=16mm]{geometry}
\usepackage{multicol,graphicx,xcolor,array,fancyhdr}
\usepackage[hidelinks]{hyperref}
\setlength{\columnsep}{6mm}
\setlength{\parindent}{1em}
\setlength{\parskip}{2pt}
\newcommand{\unk}[1]{\textcolor{red}{\texttt{\scriptsize[U+#1]}}}
\pagestyle{fancy}\fancyhf{}\renewcommand{\headrulewidth}{0.4pt}
\fancyhead[C]{\small C ARPENTRY. 631}
\begin{document}
\begin{multicols}{2}\raggedcolumns
% block 1 Text bbox=[0.0332, 0.0625, 0.0908, 0.0875]
\textbf{Currently,}

% block 2 Text bbox=[0.0898, 0.057, 0.4775, 0.1906]
As, and Ag. But because A/and-2B are equal and parallel, and Ex and f/ are also parallel; as also or and jp, it is evident, that if/s equal to eR, or to F, and sA is used to rG, or to Gg?. Therefore the forefore Ags Aso, Ap, Av, are equal to AC and AC. Therefore Ac is the compression of the beam AD, or the force pressing it on the beam HD. The proportion of these pressures, therefore, is not affected by the form of the point.

% block 3 Text bbox=[0.0918, 0.1914, 0.4785, 0.3227]
This remark is important: for many carpenters think the form and direction of the hitting joint of greater importance; and even the checks by not prosecuting the general principle through all its consequences, may be led into an order. The form of the joint is of no importance, in as far as it affects the strains in the direction of the beams, but it is often of pure consequence, in respect to its own firmness, and the effect it may have to brushing the piece on which it acts, or being crippled by it.

% block 4 Caption bbox=[0.0273, 0.3219, 0.0645, 0.3672]
\begin{center}\small\itshape Origin of the Strain on a Tbe beam.\end{center}
% block 5 Text bbox=[0.0889, 0.3234, 0.4766, 0.4281]
The same component of AB, and the same thrust on the point D by the intersection of AD, will obtain, in whatever way the original pressure on the end A is produced. Thus, supporting that a cord is made fast at all, and parallel in the direction AE, and with the same force, the beam AD will be equally compressed, and the prop D must react with the same forces.

% block 6 Text bbox=[0.0879, 0.4289, 0.4795, 0.9016]
But, it often happens to the obliquity of the pressure on AD, instead of expressing it, stretches it; and we desire to know what tension it sustains. Or they we have a familiar example in a common roof. Let the two others AC, (D) (the 18), press on the tie-beam DC. We may suppose this whole weight to B was very length between the ridge A, as if it would B were hump on there. (See Note CC). We may represent this weight by the position A of the vertical or plane lines, intersect between the ridge and the beam. Then drawing bf and agr parallel to AD and AC, Ag and Auf will represented the pressure on AC and AD. Preductor AC till CH be equal to Af. The point C is forced out in this direction, and with a force represents by the line. As this from one is not perpendicularly across the beam, it evidently stretches it; and this extending force must be withdrawn byan equal force pulling it in the opposite direction. This must arise than a similar colour thread of the opposite rather on the other end D. We concern ourselves only with this extension at present; but we see that the cohesion of the beam does nothing but supply the balance to the extending forces. It must still be supported externally, that it \_any resist\_, and by relying obligibly, be stretched. The points C and D are supported on the walls. This gives that the directions CK and DO, parallel to AA. If we draw HK parallel to DC, and HI parallel to CK (that is to Ab), meaning Dr. produced in it, it follows from the composition of forces, that the point C would be supported by the two forces KC and IC. In his manner, making DN = A, g, and completing the parallelogram DMND, the point D would be supported by the forces OD and MD. If we draw go and/k parallel to DC, it is plain that they are equal to NO and CI, while Ao and Ak are equal to

% block 7 Text bbox=[0.4932, 0.0633, 0.9287, 0.1117]
DO and CK, and Ak is equal to the sum of DO and Carpentry. CK because it is equal to Ao + A4). The weights of the roof is equal to its vertical pressure on the wall.

% block 8 Text bbox=[0.4941, 0.1125, 0.8701, 0.3914]
Thus we see, that while a pressure on A, in the direction AB, produces the strains AF and AG, on the pieces AC and AD, it has excited a strain Cl or DM in the place DC. And this completes the mechanism of a frame; for all derive their efficacy from the triangles of which they are composed, as will ape your more clearly as we proceed.

% block 9 Formula bbox=[0.5547, 0.4039, 0.7998, 0.4445]
\begin{center}\small\texttt{DO : DM : OK : DI : CM : = OO A Where Do : CK - CE: All}\end{center}
% block 10 Text bbox=[0.4912, 0.4469, 0.8623, 0.4883]
\_Or, the pressures on the points\_ C \_and D, in the direc tion of the training force A, or the increally proportioned to the portions of DE Intercepted By Al.\_

% block 11 Text bbox=[0.4893, 0.4891, 0.832, 0.5336]
Also, since A is = DO + CK, we have A\textbackslash{}; cK - CE « + d (or CD) : 50, and A\textbackslash{}; DO = DO - dz.

% block 12 Text bbox=[0.4902, 0.5367, 0.8584, 0.5906]
In general, any two of the three parallel forces AB, DO, CK, are to each other in the reciprocal proportion of the parts of OD, interconnect and their direction and the direction of the third.

% block 13 Text bbox=[0.4883, 0.5898, 0.8672, 0.782]
And this explains a still more important office of the frame aDC. I can of the grains, such as D, be supported, an external power acting at A, in the direction AB, and with an intensity which may be measured by AB, may be set in equilibrium, with another acting at C, in the direction CLs, opposite to CK or A, and with an intensity represented by CK, for since the pressure CH is partly attributed by the force IC, or the firmness of the beam DC support as D, the force K, will complete the balance. When you are not attend to the support at D, we conceive the force AB to be balanced by KC, or KC to be balanced by AK. And, in like manner, we may neglect the support or force acting at A, and consider the force DO be balanced by CK.

% block 14 Text bbox=[0.4873, 0.7828, 0.9297, 0.8594]
Thus our frame becomes a lever, and we are able it becomes to trace the interior mechanical procedure which a Leres, given it is efficacy in a by the interaction of the forces of cohesion, which cannot the points to which the external forces are applied with the supported point or flucture, and with each other.

% block 15 Text bbox=[0.4863, 0.8602, 0.8682, 0.8992]
These strains or pressure AB, DO, and CK, not being in the directions of the beams, may be called \_transverse\_. We see that by their means a frame of

\end{multicols}
\vfill\par\noindent\rule{\linewidth}{0.4pt}\par
{\scriptsize\color{gray}\noindent Source: Encyclopaedia Britannica scan, text from Parse 3.5 without correction. page\_id: britannica:309f765b1ecc83e0a7f7762e44b4a927a7555443:train:1680 \quad row index: 1680 \quad image sha256: fbb4a5fde0c2b233b6df827e5998ef2fc15cd2bc73b4eb99b93d91db83199d1b \quad parse model: nvidia/nemotron-3.5-parse \quad parse finish: stop \quad scan: \href{../scans/page-1680-fbb4a5fde0c2.jpg}{page-1680-fbb4a5fde0c2.jpg}\par}
\end{document}
